\documentclass[12pt, a4paper]{article}
\usepackage[utf8]{inputenc}
\usepackage[spanish,es-noshorthands]{babel}
\usepackage[margin=2.5cm]{geometry}
\usepackage{enumitem}

\title{\textbf{Lista de Proyectos Realizados y en Ejecución}\\Grupo de Investigación GIDEAL}
\author{Universidad de Oriente \vspace{-0.5cm}}
\date{}

\begin{document}

\maketitle

\noindent
A continuación, se detalla la relación de los principales proyectos, sistemas y herramientas desarrolladas y actualmente en ejecución por el \textbf{Grupo de Investigación en Diseño Electrónico y Aprendizaje con LLMs (GIDEAL)} dentro del espacio de trabajo multidisciplinario:

\vspace{0.5cm}

\begin{enumerate}[label=\textbf{\arabic*.}, leftmargin=*]
    \setlength\itemsep{1em}
    
    \item \textbf{Asistente Académico Inteligente (Sistema RAG):} 
    Desarrollo de un sistema conversacional de Generación Aumentada por Recuperación capaz de indexar, procesar y extraer conocimiento de material instruccional (LaTeX, Markdown, PDF) mediante bases de datos vectoriales (ChromaDB) y grandes modelos de lenguaje (vía Groq).

    \item \textbf{Agente de Diseño Electrónico Automatizado (EDA):} 
    Herramienta algorítmica capaz de parsear e interpretar esquemas de circuitos elaborados nativamente en código LaTeX (paquete \texttt{circuitikz}), extrayendo automáticamente \textit{netlists} y Listas de Materiales (BOM) para su portabilidad a softwares profesionales como EasyEDA y KiCAD.

    \item \textbf{Orquestador Multimodal para Evaluación Académica:} 
    Plataforma de automatización docente que emplea Inteligencia Artificial multimodal (visión artificial con Gemini) para analizar y hacer una evaluación preliminar (sujeta a revisión docente) de exámenes escritos manuscritos e informes de prácticas de laboratorio, generando diagnósticos estructurados y resultados directamente en informes compilables en LaTeX.

    \item \textbf{Análisis Automatizado de Diagramas y Fotografías de Circuitos:} 
    Flujo de trabajo para la ingeniería inversa y comprensión de circuitos a partir de imágenes físicas. El sistema interactúa con LLMs para identificar topologías, emitir alertas operativas y documentar todo en reportes técnicos.

    \item \textbf{Gateway de Orquestación Remota e IoT vía Mensajería:} 
    Implementación de una pasarela (Gateway) segura basada en Long Polling que actúa como interfaz entre el usuario y los agentes locales a través de Telegram, permitiendo ejecutar flujos complejos en el laboratorio sin requerir la apertura de puertos ni redes privadas.

    \item \textbf{Desarrollo de Servidores Modulares bajo el Estándar MCP:} 
    Diseño y despliegue de una arquitectura distribuida de servidores (Model Context Protocol) para evaluación, análisis de imágenes, diagnóstico del sistema y procesamiento LaTeX. Esto permite estandarizar el acceso a las herramientas del grupo.

    \item \textbf{Agentes Generadores de Material de Estudio y Evaluaciones:} 
    Creación de flujos orquestados dedicados a la elaboración de ejercicios prácticos, problemarios y propuestas de exámenes completos a demanda, alineados con el contenido de las asignaturas de Electrónica.

\end{enumerate}

\end{document}
